%% LLM-Assisted Constraint Elicitation for University Timetabling
%%
%% Target venue: ACM Transactions on Interactive Intelligent Systems.
%% Build: latexmk -pdf main.tex
%%
%% DRAFT STATUS. Sections that report measurements carry their heading and
%% nothing else: no study has been run, and a placeholder number in a results
%% section is how a draft becomes a retraction. Every quantity that will come
%% from data is wrapped in \result{}, which typesets conspicuously, so no
%% figure in this file can be mistaken for a finding.

\documentclass[acmsmall,screen,review,anonymous]{acmart}

\usepackage{booktabs}
\usepackage{array}
\usepackage{listings}

%% Conspicuous placeholder for any quantity that must come from data.
\newcommand{\result}[1]{\textbf{[\textsc{result}: #1]}}
%% A section whose content waits on data.
\newcommand{\awaitingdata}[1]{%
  \par\noindent\textit{[Awaiting data: #1]}\par}

\settopmatter{printacmref=false}
\setcopyright{none}
\acmJournal{TIIS}

\begin{document}

\title{LLM-Assisted Constraint Elicitation for University Timetabling}
\subtitle{Measuring Staff Time, Residual Error, and Trust}

\author{Anonymous Author(s)}
\affiliation{%
  \institution{Anonymous Institution}
  \country{}
}

\begin{abstract}
University timetabling constraints originate as prose: an email saying a
lecturer cannot teach on Fridays, a note that two courses must not clash, a
request for a room with workstations. Turning that prose into the formal
constraints a solver accepts is manual work, and it falls on administrative
staff. Large language models can perform the translation, and recent
benchmarks measure how faithfully they do so, but on natural language that was
generated from templates and scored without anyone in the loop. Whether the
translation helps the people who would use it is unmeasured.

We report a system in which staff write constraints in their own words, a
language model maps them onto a closed constraint vocabulary, the staff member
verifies each proposed constraint, and a solver builds the timetable from what
they accept. We evaluate it twice. An offline study measures translation
accuracy against a corpus of staff-authored constraint sentences labelled
independently by two annotators. A within-subject study with university staff
compares assisted entry against structured form entry, measuring elicitation
time with entry and verification timed separately, constraint errors remaining
after verification, and trust.

Across \result{N} staff and \result{K} labelled sentences, \result{headline
finding}. We release the labelled corpus, the prompts, the code, and the
pre-registration.
\end{abstract}

\begin{CCSXML}
<ccs2012>
<concept>
<concept_id>10003120.10003121.10003129</concept_id>
<concept_desc>Human-centered computing~Interactive systems and tools</concept_desc>
<concept_significance>500</concept_significance>
</concept>
<concept>
<concept_id>10010147.10010178.10010187</concept_id>
<concept_desc>Computing methodologies~Knowledge representation and reasoning</concept_desc>
<concept_significance>300</concept_significance>
</concept>
<concept>
<concept_id>10003456.10003457.10003527</concept_id>
<concept_desc>Social and professional topics~Computing education</concept_desc>
<concept_significance>100</concept_significance>
</concept>
</ccs2012>
\end{CCSXML}

\ccsdesc[500]{Human-centered computing~Interactive systems and tools}
\ccsdesc[300]{Computing methodologies~Knowledge representation and reasoning}
\ccsdesc[100]{Social and professional topics~Computing education}

\keywords{constraint elicitation, timetabling, large language models,
human-in-the-loop verification, constraint programming, user study}

\maketitle

%% =====================================================================
\section{Introduction}
%% =====================================================================

A university timetable is the output of a constraint satisfaction problem, and
most of its constraints begin as sentences. A lecturer writes that she cannot
teach on Fridays. A programme coordinator notes that two courses fell in the
same slot last year and that students could not attend both. A technician
records that a practical needs thirty-eight workstations. Somebody then reads
these and writes them into whatever form the scheduling system accepts, and
that somebody is usually an administrator with no background in constraint
programming.

The translation is the bottleneck, and it is a human one. Solver technology for
curriculum-based timetabling is mature: the ITC-2007 instances have been
attacked for close to two decades~\cite{digaspero2007itc}, and a systematic
review of metaheuristics for course, exam and school timetabling covers the
last three years alone. What remains expensive is getting the problem stated.

Large language models are an obvious candidate, and the constraint-programming
community has taken up the idea. CP-Bench measures how well models produce
constraint models from problem descriptions and names modelling as the
bottleneck for wider adoption of constraint
programming~\cite{michailidis2025cpbench}. ConstraintLLM pairs a fine-tuned
model with solver-checked self-correction~\cite{shi2025constraintllm}.
SCHEDBench puts scheduling problems, curriculum timetabling among them, to
thirteen models and finds that none is invariant to how the same problem is
worded~\cite{sharma2026schedbench}. All three measure the model. None measures
what happens when a person is the one supplying the words and checking the
output.

That distinction matters more than it may appear, for two reasons. The first is
the input. SCHEDBench renders its natural language with a deterministic
template pipeline, and its own limitations section names broader
human-authored phrasing as the route to better ecological validity. A model
that handles templated prose need not handle an email from a lecturer. The
second is the loop. In any deployment a person has to confirm that the
formalised constraint says what they meant, and verification is not free: the
closest timed comparison of a natural-language interface against a manual
alternative found the interface 31\% faster per query and no more accurate,
with correct output per hour unchanged, because the time saved on composing
reappeared on checking~\cite{ipeirotis2026nlidb}. A system evaluated only on
translation accuracy cannot see that effect at all.

This paper measures it. We built a system in which staff write constraints in
their own words; a language model maps each sentence onto a closed vocabulary
of constraint types; a resolver binds the entities named in the sentence to
the department's own courses, rooms, lecturers and groups; rule-based checks
flag what needs attention; the staff member accepts, edits or rejects each
proposed constraint; and a metaheuristic solver builds a timetable from the
accepted set. The language model never produces or edits a timetable.

Our contributions are:

\begin{itemize}
  \item A corpus of staff-authored timetabling constraint sentences, labelled
        independently by two annotators with inter-annotator agreement
        reported. Existing scheduling benchmarks template their natural
        language and report no agreement statistic.
  \item A measurement of translation accuracy on that corpus, decomposed into
        type, strength and entity accuracy, with entity errors separated into
        references the system could not resolve and references it resolved to
        the wrong entity.
  \item A within-subject study with university staff comparing assisted entry
        against structured form entry, with elicitation time reported as entry
        time and verification time separately, constraint errors remaining
        after verification as a pre-registered non-inferiority check, and
        trust and usability as secondary measures.
  \item The system, the prompts, the corpus and the pre-registration, released
        for reuse.
\end{itemize}

Section~\ref{sec:study1} reports the offline measurement and
Section~\ref{sec:study2} the staff study.

%% =====================================================================
\section{Related Work}
%% =====================================================================

\subsection{Natural Language to Formal Models}

The line from NL4Opt~\cite{ramamonjison2022nl4opt} through CP-Bench
\cite{michailidis2025cpbench} and ConstraintLLM~\cite{shi2025constraintllm}
treats the problem as a translation task scored against reference models. The
benchmarks have grown more demanding, and later work documents annotation
errors in the earlier ones, which is part of why we label our corpus twice and
report agreement.

SCHEDBench~\cite{sharma2026schedbench} is the closest of these to timetabling.
It covers six scheduling families including curriculum timetabling on the
ITC-2007 Track 3 instances, and reports that feasibility collapses under
surface-form variation. Two properties bound what it can tell us. Its natural
language comes from a deterministic template pipeline with slot filling, so no
domain user wrote any of it; and it scores end-to-end schedule generation,
stating that its invariance claims do not isolate effects to the modelling
stage. Its poor feasibility results are, if anything, an argument for the
architecture we adopt, in which the model emits constraints and never a
timetable.

\subsection{Interactive Constraint Acquisition}

Constraint acquisition asks a user questions and learns a constraint network
from the answers~\cite{bessiere2023partialqueries}. LlmAcq places a language
model inside that loop, so a user may answer in natural language, and reports a
large reduction in the number of queries needed~\cite{mechqrane2026llmacq}.

Two features of its evaluation bear on ours. The domains are logic puzzles, not
scheduling. And the human element is a text-collection exercise: twelve
students with no constraint-programming background wrote one sentence per
Sudoku grid, and those 120 sentences were replayed in simulation. No time,
error rate or trust measurement is reported. Its two substantive findings about
human feedback shaped our design, and we return to them in
Section~\ref{sec:resolution}: people do not answer as specifically as an
algorithm would prefer, and they rely on implicit structural knowledge, so a
sentence may name no variable at all.

\subsection{Human-in-the-Loop Verification}

MeetMate~\cite{lawless2024meetmate} is the nearest system to ours. A chat
interface collects meeting constraints, a model decides whether each can be
coded, and a code-generating component emits a predicate per constraint;
candidate times are enumerated and scored by weighted sum. Its evaluation is a
single-condition technology probe with ten information workers at one company.
The paper declines to report quantitative outcomes it collected, on the
grounds that the open-ended conversational design makes them invalid, and it
reports no baseline, no time on task, and no usability or trust instrument. It
has no hard constraint notion, treating everything as a weighted preference,
and no constraint-level review step: participants never saw the constraint
list, which the authors identify as why they could not troubleshoot the
system's errors.

VeriPlan~\cite{lee2025veriplan} is the nearest verification interface. Prompt
constraints are translated into temporal-logic rules, each rule is presented
back in natural language for the user to confirm, a per-rule slider sets
enforcement strictness, and a model checker verifies generated plans against
the confirmed rules. Its four-condition within-subject ablation, including an
arm with no checker at all, is the design our optional third condition adapts.
Its measures are perceived usefulness, ease of use and satisfaction; it
reports no task time, no error count and no trust scale, and its domain is
personal and everyday planning.

OptiChat~\cite{chen2026optichat} explains optimisation models to non-experts
and reports descriptions and infeasibility diagnoses in under a minute against
eighteen to fifty-five minutes for human experts. The comparison is against
expert effort, not against the effort of the staff who supply the constraints.
TRACE-CS~\cite{vasileiou2025tracecs} answers contrastive questions about
course schedules and, like our system, confirms the parsed query with the user
before acting; its explanation correctness was graded by its authors.

\subsection{Timetabling and Education Administration}

The timetabling literature is dominated by solution methods. Two 2026
systematic reviews of large language models in education cover teacher
workload and instructional planning~\cite{feduc2026review} and intelligent
education systems broadly~\cite{info2026review}; neither reaches timetabling.
The administrative side has attitude surveys without systems, while the
constraint-programming and HCI literatures have systems without institutional
users.

\subsection{Position of This Work}

Table~\ref{tab:position} sets out the elements this study combines against the
nearest work on each. Every element has a near neighbour. The combination does
not, and it is the combination that the measurement requires: a corpus written
by the people the system is for, labelled with agreement; a baseline that is
the strongest realistic alternative; time split into composing and checking;
and error measured after verification, which is the only point at which the
question ``did this help?'' can be asked.

\begin{table}[t]
  \caption{Elements of this study against the nearest prior work on each.}
  \label{tab:position}
  \small
  \begin{tabular}{@{}p{0.27\linewidth}p{0.36\linewidth}p{0.29\linewidth}@{}}
    \toprule
    \textbf{Element} & \textbf{Nearest prior work} & \textbf{This study} \\
    \midrule
    Human-authored constraint text
      & SCHEDBench templates its natural language~\cite{sharma2026schedbench}
      & Corpus written by staff \\
    Inter-annotator agreement
      & Not reported~\cite{lawless2024meetmate,sharma2026schedbench}
      & Cohen's $\kappa$ on type and strength~\cite{cohen1960kappa} \\
    Non-LLM baseline
      & Single-condition probe~\cite{lawless2024meetmate}
      & Structured form, identical for all \\
    Measured time
      & Perceived efficiency~\cite{lee2025veriplan}
      & Entry and verification timed apart \\
    Error after verification
      & Not measured
      & Pre-registered non-inferiority check \\
    Trust instrument
      & Dropped as unreliable~\cite{ipeirotis2026nlidb}
      & Named in the pre-registration, piloted \\
    \bottomrule
  \end{tabular}
\end{table}

%% =====================================================================
\section{System}
%% =====================================================================

Five components: a data layer holding the department's courses, rooms, slots,
lecturers and student groups; a constraint schema; a translation component; a
verification interface; and a solver. Staff text enters at the translation
component and leaves as constraints. The model is never given a timetable and
never asked to produce one.

\subsection{Constraint Vocabulary}

The vocabulary is closed and small, because the translation target has to be
something a model can hit reliably and a person can read back. Eight types are
solvable: lecturer unavailability, lecturer time preference, a course pinned
to a named room, a course requiring a type of room, a cap on consecutive
periods, an earliest acceptable period, a pair of courses that must not
overlap, and a set of courses to be placed on one day. Every constraint
carries a strength, hard or soft, a soft weight on a one-to-ten scale, and the
span of source text it came from.

Two further types stay open and never reach the solver.
\texttt{unsupported} marks a schedulable request the vocabulary cannot
express; \texttt{needs\_clarification} marks input too ambiguous to formalise
and carries the question to put to the user. Both are logged and categorised.
Their distribution is a finding about what staff ask for, and we report it as
one, not as a defect rate.

\subsection{Translation}

A system prompt carries the vocabulary, the department's register of entities,
and sixteen worked examples. The model is asked for JSON, the reply is
validated against the schema, and a single retry is issued on a validation
failure with the validation error appended. Three outcomes are never accepted
as a result: a refusal, a response truncated at the token ceiling, and an API
error. Each is recorded as that item's outcome and counted.

Provider choice is a parameter. Ten providers sit behind one interface, and a
run is identified by a provider and model pair. Two provider features that
would break symmetry are disabled by default: native JSON mode, which the
reference configuration does not use, and sampling temperature, which the
models used here do not accept. An instrument that differs between arms
measures the harness.

\subsection{Entity Resolution}
\label{sec:resolution}

Staff name entities the way people do. The model copies their wording into the
entity fields, and resolution against the department register happens
afterwards, which keeps the output contract simple and makes resolution
failures measurable on their own.

Resolution records how each reference was matched: by identifier, by full
name, by a recorded alias, or by token containment. It also records whether
the reference was \emph{descriptive}, which we define as a reference that
picks out an entity without naming it: ``the big lab'', ``the first years'',
``my Tuesday group''. This class is the one LlmAcq found people actually
produce~\cite{mechqrane2026llmacq}, and it fails differently from a
hallucinated name: a description resolved to the wrong room is a resolution
error, while an invented lecturer is a fabrication. We count them separately,
and report accuracy on descriptive references as its own figure.

\subsection{Flagging}

Flagging is by rule. Self-reported model confidence is not used: it is
uncalibrated, and the field is present on the schema but neither requested nor
read, so a calibration study remains a configuration change and not a schema
change.

A blocking flag keeps a constraint away from the solver until a person
resolves it: an unresolved entity, an ambiguous entity, an open category, or a
period outside the configured grid. A review flag surfaces the constraint
prominently and leaves the decision to the user: a descriptive reference, a
weak partial match, a strength cue contradicting the parsed strength, or a
room too small for its course. Verification is mandatory under either.

\subsection{Verification Interface}

Each proposed constraint is shown as a sentence beside the words it came from
--- ``Dr.\ Sarah Lee cannot teach on Friday: \textsc{hard}'' --- with accept,
edit and reject controls and a hard/soft toggle. Raw JSON is never shown. A
participant who cannot read the restatement is not verifying anything, and the
residual-error measure would then be measuring the interface.

The interface logs every action with a timestamp, and derives durations from
paired phase spans. Entry time and verification time are therefore separate
quantities, recorded and reported apart. This is a direct response to
\cite{ipeirotis2026nlidb}, which bundled composing and checking into one
completion time and could not separate them.

\subsection{Solver}

Timetables are assignments of courses to a room and a starting slot. Cost is a
large penalty for each structural or hard elicited violation plus the weighted
sum of soft violations, with the hard penalty set so that no accumulation of
soft penalties can purchase a hard breach. Each constraint type contributes a
violation count, so a timetable breaching one constraint twice scores worse
than one breaching it once. Search is simulated annealing over three move
types, seeded and deterministic. The solver is infrastructure here, not a
contribution; we describe it so the measurements are reproducible.

When the accepted constraints cannot all hold, the solver reports which hard
constraints were breached, identified by the staff sentence they came from, so
the conflict can be put back to the person who stated it.

%% =====================================================================
\section{Study 1: Translation Accuracy}
\label{sec:study1}
%% =====================================================================

\subsection{Corpus}

Between 150 and 300 constraint sentences written by university teaching and
administrative staff. The requirement that they be human-authored is the point
of the corpus: existing scheduling benchmarks template their natural language,
and SCHEDBench names human-authored phrasing as what it
lacks~\cite{sharma2026schedbench}.

\subsection{Labelling}

Two annotators independently label each sentence with the constraint or
constraints it expresses, using canonical entity identifiers. Agreement is
reported as Cohen's $\kappa$~\cite{cohen1960kappa} on the type assignment and
on the strength assignment, with the unit of analysis stated: $\kappa$ requires
one category per item, so the primary constraint of each sentence supplies the
label, and the number of sentences carrying more than one constraint is
reported alongside so the cost of that convention is visible. Set-level exact
agreement is reported using every constraint. Disagreements are adjudicated to
a gold label.

\subsection{Measures}

Exact match at the sentence level; type accuracy; hard/soft accuracy
conditional on a correct type; and entity accuracy, decomposed into correct,
unresolved, misresolved, missed and spurious, with descriptive references
counted separately.

Three comparison conventions are fixed in advance. Source-text spans are
excluded, since annotators disagree about boundaries without disagreeing about
the reading. Free-prose fields on the two open categories are excluded for the
same reason. Soft weights are excluded by default, a weight being a scale
choice and not a reading of the sentence, with a weight-inclusive variant
reported as a sensitivity analysis.

\subsection{Sensitivity}

Prompt wording, model choice and reasoning effort are varied, with every run
identified by its model, prompt digest and effort level, and every raw
response recorded so a scored run can be re-scored without further calls.

%% =====================================================================
\section{Study 2: Staff Study}
\label{sec:study2}
%% =====================================================================

\subsection{Design}

Within-subject, with order counterbalanced across participants. Target sample
is 25 to 30 university staff who hold timetabling responsibility. The analysis
plan, the primary outcome and the success thresholds are pre-registered before
data collection.

\paragraph{Condition 1 (baseline): structured form entry.}
Dropdowns for lecturer, day, period, room and strength, identical for every
participant. This is the strongest realistic non-LLM alternative; a comparison
against an unstructured email baseline would support a weaker claim. Nothing
in this condition can be misread, so it has no verification phase.

\paragraph{Condition 2: assisted entry.}
Free text through the translation component, followed by constraint-level
verification.

\paragraph{Condition 3 (optional): assisted entry without verification.}
Adapted from VeriPlan's no-checker arm~\cite{lee2025veriplan}, this separates
the contribution of the verification interface from that of the translation
component. It is included if the within-subject burden at the target sample
allows.

\subsection{Tasks}

Each participant expresses a scenario given as a task brief. Briefs are
presented as shaded week grids, room and equipment inventories, complaint logs
and email extracts, never as the sentences a participant should produce: a
brief that states the target sentence would make the study a transcription
race, in which the free-text arm copies and the form arm transcribes.

Each task carries between six and ten intended constraints. Smaller tasks make
the residual-error measure too coarse for the margin it has to test: with
three intended constraints a participant's error rate can only take the values
$0$, $0.33$, $0.67$ and $1$, a granularity larger than the non-inferiority
margin itself.

\subsection{Outcomes}

\paragraph{Primary.} Elicitation time per constraint set, assisted against
baseline, with a target reduction fixed in the pre-registration. Entry time
and verification time are reported separately and summed to the total.

\paragraph{Secondary: residual error.} Constraint errors remaining after
verification, scored against the task key, treated as a non-inferiority check
with an absolute percentage-point margin. The gap between the set a
participant finally accepted and the set the task asked for decomposes into
missed, spurious, wrong type, wrong strength, wrong entity and unresolved.

\paragraph{Secondary: throughput.} Correct constraints per minute of staff
time. Neither of the outcomes above detects a time reduction cancelled by an
accuracy loss; \cite{ipeirotis2026nlidb} found exactly that pattern, with
correct output per hour unchanged.

\paragraph{Secondary.} Usability by the System Usability
Scale~\cite{brooke1996sus}, a trust scale named in the pre-registration,
perceived acceptability of the resulting timetable, and a short interview.

\subsection{Joint Success Rule}

The claim holds only if the primary target is met and the non-inferiority
margin is not breached. A pilot sets both thresholds, which are locked before
pre-registration and reported as absolute margins in percentage points.

\subsection{Analysis}

Paired $t$-test or Wilcoxon signed-rank according to the pre-registered
decision rule, with effect sizes and confidence intervals. Sessions that were
abandoned or left with an open phase are excluded from the timing analysis and
reported as exclusions, since an unclosed phase is closed at the last recorded
event and would otherwise distort a condition mean without bound.

%% =====================================================================
\section{Results}
\label{sec:results}
%% =====================================================================
\awaitingdata{Neither study has been run; ethics approval for Study~2 is
pending. The subsection headings are retained to fix the reporting order.}

\subsection{Translation Accuracy}
\subsection{Error Categories}
\subsection{Cross-Model Comparison}
\subsection{Elicitation Time}
\subsection{Residual Error}
\subsection{Throughput}
\subsection{Trust and Usability}

%% =====================================================================
\section{Discussion}
%% =====================================================================
\awaitingdata{Depends on Section~\ref{sec:results}.}

%% =====================================================================
\section{Threats to Validity}
%% =====================================================================

Several threats are structural and can be stated before any data exists.

\paragraph{The baseline is a choice.} A structured form is the strongest
realistic non-LLM alternative, but it is not the only one. Prior work finds
that a model suggesting values \emph{into} a structured form cuts completion
time substantially~\cite{acm2025formcompletion}, and that hybrid is neither of
our main conditions. We compare free text against a form because that is the
comparison our research questions ask, and we report the omission.

\paragraph{Verification time is attributable only under a split.} Reporting a
single elicitation total would leave the central objection untestable. We
report entry and verification separately for that reason, but the split
depends on participants moving between interface phases in the way the
instrument assumes, and a participant who verifies while still composing is
mismeasured.

\paragraph{Timing is sensitive to deployment.} The primary outcome is a
duration measured through a browser interface. Remote hosting introduces
network latency and its variance into the dependent variable, so sessions are
run locally.

\paragraph{Task keys encode a scoring decision.} Residual error is scored
against what the task asked for, which means the researcher's reading of the
brief is the standard. Where a brief admits a defensible alternative
formalisation, the decision is recorded in the pre-registration, not settled
once the data arrive.

\paragraph{Open categories are not errors.} A participant who recognises that
the vocabulary cannot express a request has answered correctly. Treating
\texttt{unsupported} as a failure would reward forcing requests into
near-fitting types.

\paragraph{Generality.} One department, one constraint vocabulary, one
timetabling formulation. The vocabulary was frozen early to keep the
translation target stable, which bounds what staff could express and is itself
measured through the open categories.

%% =====================================================================
\section{Conclusion}
%% =====================================================================
\awaitingdata{Depends on Section~\ref{sec:results}.}

%% =====================================================================
\section*{Availability}
%% =====================================================================

The system, the prompts with their digests, the labelled corpus with both
annotators' independent labels and the adjudicated gold, the raw model
responses, the task keys, and the pre-registration are released. The raw
responses allow every accuracy figure to be re-scored without access to a
model provider.

\bibliographystyle{ACM-Reference-Format}
\bibliography{references}

\end{document}
